\chapter{Lay Summary}
Modern computer chips rely on internal networks to move data between many processing units, but these connections can wear out over time, especially when some paths are used more than others. Traditional routing methods send data along fixed paths, which can overload certain parts of the network and shorten the chip’s lifespan. This thesis explores a learning-based approach where the network adapts its routing decisions based on current conditions, spreading traffic more evenly to reduce stress on heavily used components. By combining this approach with models of how hardware wears out, the system learns to avoid paths that are likely to fail sooner. Our results show that this method can significantly extend the lifetime of the network while maintaining or even improving communication speed, highlighting a promising direction for building more reliable and efficient computer systems.